\documentclass[11pt]{article}
\usepackage[margin=1in]{geometry}
\usepackage{amsmath,amssymb}
\usepackage{enumitem}
\title{Regime Canonicalization and Provenance Refactor in \texttt{BindingAndCatalysis.jl}}
\author{Internal Technical Report}
\date{\today}
\begin{document}
\maketitle

\section{Problem Statement}
The prior implementation primarily stored per-regime dense/sparse matrices $(C_x,C_{0,x})$, $(H,H_0)$, and $(C_{qK},C_{0,qK})$ directly inside each vertex object. This duplicated geometric information and obscured causal provenance. In particular, the representation did not expose the chain
\[
\text{x-space boundary atom} \longrightarrow \text{mapping object} \longrightarrow \text{qK-space hyperplane atom}.
\]

\section{New Architecture}
The refactor introduces factorized pools and reference-based ownership.
\subsection*{A. Canonical x-space atoms}
Each x-space boundary is stored once as an \texttt{XHalfspaceAtom} keyed by
\((i,\min(j_1,j_2),\max(j_1,j_2))\).
Orientation is represented via \texttt{XHalfspaceRef(atom\_id, sign)}.

\subsection*{B. Mapping-object pools}
Regular regimes pool \texttt{RegularMapping(H,H0)} objects.
Singular regimes pool \texttt{SingularMapping(nullity, H\_like, right\_null, left\_null)}.
This makes singular regimes first-class and preserves geometry-relevant null-space data.

\subsection*{C. qK hyperplane atoms with provenance}
A qK inequality/hyperplane is pooled as \texttt{QKHyperplaneAtom(a,b,provenance)}, where provenance is a set-like list of pairs \((x\_atom\_id, mapping\_id)\).
Thus multiple causal sources may map to the same qK hyperplane atom without erasing lineage.

\subsection*{D. Lightweight regime ownership}
Each vertex now stores:
\begin{itemize}[leftmargin=2em]
\item \texttt{x\_halfspace\_refs}
\item \texttt{mapping\_id}
\item \texttt{qk\_atom\_ids}
\end{itemize}
and can still materialize $(C_x,C_{0,x})$ and $(C_{qK},C_{0,qK})$ for compatibility.

\subsection*{E. Regime-graph edge provenance}
Each \texttt{VertexEdge} now carries:
\begin{itemize}[leftmargin=2em]
\item \texttt{x\_atom\_id}, \texttt{x\_atom\_sign}
\item optional \texttt{qk\_atom\_id}
\end{itemize}
so edge transitions are geometrically interpretable.

\section{Design Rationale}
\begin{enumerate}[leftmargin=2em]
\item \textbf{Canonicalization before materialization}: removes structural duplication early.
\item \textbf{Direction separate from hyperplane identity}: avoids atom duplication for sign flips.
\item \textbf{Singular mapping as explicit type}: prevents opaque handling of nullity-1 regimes.
\item \textbf{Provenance-preserving qK pooling}: supports interface interpretation and overlap/layering analysis.
\item \textbf{Graph-edge metadata}: prepares for direct edge-based propagation algorithms.
\end{enumerate}

\section{Rank-1 Update Compatibility}
The storage now decouples regime identity from fully materialized matrices. Neighboring regimes differing by one dominance choice can reuse edge metadata and mapping references, creating a natural insertion point for future Sherman--Morrison style updates for regular mappings and corresponding singular-branch transitions.

\section{API Additions}
New queries were added to inspect and debug the factorized model:
\begin{itemize}[leftmargin=2em]
\item regime-level: \texttt{get\_x\_halfspace\_refs}, \texttt{get\_mapping\_id}, \texttt{get\_qk\_atom\_ids}
\item global pools: \texttt{get\_x\_halfspace\_atoms}, \texttt{get\_regular\_mappings}, \texttt{get\_singular\_mappings}, \texttt{get\_qk\_hyperplane\_atoms}
\item edge-level: \texttt{get\_edge\_provenance}
\end{itemize}

\section{Limitations and Future Work}
This pass establishes data-model support; it does not yet implement the full edge-walk incremental solver. A next phase should:
\begin{itemize}[leftmargin=2em]
\item implement rank-1 updates for \texttt{RegularMapping}
\item derive nullity-1 transitions from previous inverse data directly
\item define deterministic row correspondence for singular qK constraints where elimination reorders constraints.
\end{itemize}

\end{document}
